% Internal proposal: this file is intentionally excluded from submission builds.
\documentclass[sigconf,anonymous]{aamas}
% Copyright and conference block copied from the official AAMAS 2027 sample.
\makeatletter
\gdef\@copyrightpermission{
  \begin{minipage}{0.2\columnwidth}
   \href{https://creativecommons.org/licenses/by/4.0/}{\includegraphics[width=0.90\textwidth]{by}}
  \end{minipage}\hfill
  \begin{minipage}{0.8\columnwidth}
   \href{https://creativecommons.org/licenses/by/4.0/}{This work is licensed under a Creative Commons Attribution International 4.0 License.}
  \end{minipage}
  \vspace{5pt}
}
\makeatother
\setcopyright{ifaamas}
\acmConference[AAMAS 2027]{Proc.\@ of the 26th International Conference
on Autonomous Agents and Multiagent Systems (AAMAS 2027)}{3--7 May 2027}{Hanoi, Vietnam}{M.~Baldoni, F.~Fang, W.~Yeoh, N.~Yorke-Smith (eds.)}
\copyrightyear{2027}
\acmYear{2027}
\acmDOI{}
\acmISBN{}
\acmSubmissionID{INTERNAL REVISION}
\submissionType{Research Paper Track}

\acmSubmissionID{INTERNAL PRE-RESULTS PROPOSAL}

\title[Earning Roles]{Earning Roles: Self-Evolving Responsibility from Situated Peer Judgments}
\ccsdesc[500]{Computing methodologies~Multi-agent systems}
\keywords{Multiagent coordination, self-evolving roles, peer judgment, online adaptation}

\begin{abstract}
Multi-agent teams must assign responsibility before a dependent task is executed,
yet static roles, aggregate reputation, and terminal rewards do not record whether
a particular delivery helped the agent that used it. We introduce Earning Roles, a
framework in which a recipient's situated judgment---grounded in its use,
modification, rejection, or responsibility-checked repair of a versioned
delivery---becomes attributable evidence for the producer's future
responsibilities. Evidence arrives after execution and may be revised by later
independent outcomes; a subsequent owner can therefore use the evolving role state
before choosing a peer. We formalize this process under selected-only feedback,
delayed updates, local candidate sets, and explicit producer, recipient, and
downstream responsibility boundaries. We evaluate whether the resulting assignments
improve unseen-task quality and complete collaboration cost, using genuine
producer--recipient dependencies, held-out task roots, independent checks, and
matched model, information, exploration, and communication budgets. Comparisons
with raw acceptance, contextual trust and bandit selection, terminal-only feedback,
and pooled-history control isolate the value of attributable situated judgment. The
framework is designed to support online updates while measuring adaptation to drift,
update latency, and retention of prior competence.
\end{abstract}

\begin{document}
\maketitle

\noindent\textbf{Internal pre-results proposal --- not for submission.}
This document fixes a working question and prospective tests. It does not
report an evaluated learning method. The benchmark, backbone, update mechanism,
and empirical conclusions remain open.

\section{The Contribution We Need to Establish}

The target is a mechanism that uses situated judgments to change useful future
responsibilities. It must specify who may observe and share an evaluation,
what can be updated after delayed evidence arrives, and how the resulting
state changes a decision before work begins. A role label by itself is not an
outcome. The intended evidence is a change in actual assignments that improves
subsequent quality and complete collaboration cost.

The first development task separates a producer's code delivery from a
recipient's dependent integration work. These are temporary responsibilities,
not preset expert identities. A short, single-handoff probe can establish an
executable information boundary; it cannot establish decentralized role
formation or the full recursive collaboration setting.

\section{Prospective Discriminating Tests}

\paragraph{Task and signal qualification.}
Use real delivered work and record the recipient's initial judgment, actual
use, revisions, and completion cost. Keep independent evaluation outside the
agent observation boundary. Distinguish the quality of a delivery from the
quality obtained after another agent repairs it. Split by task structure
rather than treating renamed instances as independent problems.

\paragraph{Future responsibility.}
Evaluate legal pre-execution assignments on unseen tasks, including decisions
by collaborators who did not directly judge the candidate. Match initial
agents, observations, exploration opportunities, memory, and total resources.
Give the contextual trust or bandit comparator the same evidence. Test whether
task-dependent divisions of responsibility explain any benefit beyond selecting
one globally stronger peer.

\paragraph{Mechanism and cost.}
Remove judgment information, scramble its source within comparable contexts,
or withhold its transfer while preserving the remaining execution rules. If
real-time training becomes the contribution, compare one explicit updater with
online SGD, RLS, and periodic refitting. Measure old-task retention, adaptation,
and feedback-to-serving latency jointly; inexpensive parameter updates alone
do not establish useful real-time learning.

\section{Current Evidence Boundary}

The existing executable fixtures test protocol ordering and artifact lineage.
Their judgments and repairs are programmed, so they do not support learning
efficacy. Earlier API probes establish narrower transport and interaction
observations. There is no completed controlled result for the method proposed
here. All result claims and a final algorithm description must be supplied from
auditable experiments before this proposal can become a submission manuscript.

\end{document}
